\section{基本函数的导数}
\label{sec:ru-elementary-derivatives}
\label{sec:analysis-elementary-derivatives}

有了求导规则，复杂函数的求导就归结为少数基本函数。本节先由有限运算和上一章的基本极限得到幂函数、指数和三角函数的导数，
再借反函数与复合的规则得到其余基本函数的导数，并逐项注明成立范围。

常数函数的差商恒为零，恒等函数的差商恒为$1$。对正整数$m$，把$x^m$看成$m$个恒等函数之积，反复使用乘积法则，每一项只对一个因子求导，共$m$项，得$(x^m)'=mx^{m-1}$。
对负整数$m=-n$，在$x\neq0$处由商法则得$(x^{-n})'=-nx^{n-1}/x^{2n}=mx^{m-1}$，形式相同。于是多项式可以逐项求导，例如$(x^3-3x)'=3x^2-3$。

指数函数满足$e^{x+h}=e^xe^h$，由第~\ref{sec:ru-elementary-limits}节的基本极限$(e^h-1)/h\to1$，
\[
 \frac{e^{x+h}-e^x}h=e^x\,\frac{e^h-1}h\to e^x,\qquad\text{即}\ (e^x)'=e^x.
\]
正弦和余弦由加法公式把差商写成两个基本极限的组合：
\[
 \frac{\sin(x+h)-\sin x}h=\sin x\,\frac{\cos h-1}h+\cos x\,\frac{\sin h}h\to\cos x,
\]
同理$(\cos x)'=-\sin x$。这两个推导只用到加法公式和基本极限，不依赖尚未证明的逐项求导；第~\ref{sec:ru-termwise-differentiation}节对指数级数逐项求导会得到同样的结果，可作核对。
由商法则，$(\tan x)'=(\cos^2x+\sin^2x)/\cos^2x=1/\cos^2x=1+\tan^2x$，在$\cos x\neq0$处成立。

其余基本函数是这些函数的反函数或复合。自然对数是指数函数的反函数，指数的导数处处非零，由反函数求导，$(\log x)'=1/e^{\log x}=1/x$（$x>0$）。
反正弦是正弦在$[-\pi/2,\pi/2]$上的反函数。在$(-\pi/2,\pi/2)$内$(\sin)'=\cos>0$，所以对$y\in(-1,1)$，
\[
 (\arcsin y)'=\frac1{\cos(\arcsin y)}=\frac1{\sqrt{1-y^2}},
\]
其中$\cos(\arcsin y)\geq0$保证开方取正号。在$y=\pm1$处，对应的$\cos=0$，反正弦不可导，图像有竖直的切线。
同理$(\arccos y)'=-1/\sqrt{1-y^2}$（$-1<y<1$），而$(\arctan y)'=1/(1+\tan^2(\arctan y))=1/(1+y^2)$对所有实数$y$成立。

一般底数与实数次幂由自然指数和对数定义：$a^x=e^{x\log a}$（$a>0$），$x^\alpha=e^{\alpha\log x}$（$x>0$）。由链式法则，
\[
 (a^x)'=a^x\log a,\qquad (x^\alpha)'=e^{\alpha\log x}\cdot\frac\alpha x=\alpha x^{\alpha-1},\qquad (\log_ax)'=\frac1{x\log a}.
\]
$x^\alpha$的公式只在$x>0$上使用，因为一般实数次幂只对正数定义；$\alpha$为整数时它与前面的结果一致，并且后者在负数处也成立。
表~\ref{tab:analysis-elementary-derivatives}汇总这些公式。把它们与四则运算、链式法则和反函数求导结合，就能求出基本函数经有限次运算与复合得到的函数的导数。

\begin{table}[htbp]
 \centering\small
 \begin{tabular}{@{}lll@{}}
 \toprule
 函数 & 导数 & 成立范围\\
 \midrule
 $x^m$（$m\in\mathbb Z$） & $mx^{m-1}$ & $m\geq0$时$x\in\mathbb R$；$m<0$时$x\neq0$\\
 $x^\alpha$（$\alpha\in\mathbb R$） & $\alpha x^{\alpha-1}$ & $x>0$\\
 $e^x$ & $e^x$ & $x\in\mathbb R$\\
 $a^x$（$a>0$） & $a^x\log a$ & $x\in\mathbb R$\\
 $\log x$ & $1/x$ & $x>0$\\
 $\log_ax$（$a>0$，$a\neq1$） & $1/(x\log a)$ & $x>0$\\
 $\sin x$ & $\cos x$ & $x\in\mathbb R$\\
 $\cos x$ & $-\sin x$ & $x\in\mathbb R$\\
 $\tan x$ & $1/\cos^2x$ & $x\neq\pi/2+k\pi$\\
 $\arcsin x$ & $1/\sqrt{1-x^2}$ & $-1<x<1$\\
 $\arccos x$ & $-1/\sqrt{1-x^2}$ & $-1<x<1$\\
 $\arctan x$ & $1/(1+x^2)$ & $x\in\mathbb R$\\
 \bottomrule
 \end{tabular}
 \caption{基本函数的导数及成立范围。三角函数采用弧度制，$\log$为自然对数。}
 \label{tab:analysis-elementary-derivatives}
\end{table}
